%% kapa: exam=1 round=29 year=2018 period=1 form=A subject=civil_law subject_ko=민법 seq=17 number=17
%% source: Q-Net 감정평가사 기출문제 (cst00309) artlSeq=5202808 file=extracted/1008/5202808/2018년 제29회 감정평가사 1차시험 기출문제/2018년 제29회 감정평가사 1차시험 기출문제/1교시 A형.pdf
%% conversion: kordoc 4.24.51 -> tools/gen1.py (auto); review=proofread-sonnet; flags=none
\begin{question}{17}
소멸시효 완성에 관한 설명으로 옳지 않은 것은? (다툼이 있으면 판례에 따름)
\begin{choices}
  \item 소유권은 소멸시효에 걸리지 않는다.
  \item 동일한 목적을 달성하기 위하여 복수의 채권을 가진 채권자가 어느 하나의 채권만을 행사하는 것이 명백한 경우, 채무자의 소멸시효 완성 항변은 채권자가 행사하는 당해 채권에 대한 항변으로 볼 수 있다.
  \item 유치권이 성립된 부동산의 매수인은 피담보채권의 소멸시효 완성으로 직접 이익을 받는 자에 해당하지 않으므로 소멸시효의 완성을 원용할 수 없다.
  \item 물상보증인은 피담보채권의 소멸에 의하여 직접 이익을 받는 관계에 있으므로 피담보채권의 소멸시효의 완성을 주장할 수 있다.
  \item 채무불이행으로 인한 손해배상청구권에 대한 소멸시효 항변이 불법행위로 인한 손해배상청구권에 대한 소멸시효 항변을 포함한 것으로 볼 수는 없다.
\end{choices}
\end{question}
\answer{3}
